% GENERATED FILE. Edit canonical lessons, then run npm run generate:books.
% canonical-source-hash: fnv1a64:1d743d513a531d7d
% canonical-lessons: ES-C44-plural-es

\chapter{-es --- The Plural Chapter One Promised}
\label{ch:plural-es}
\hlchaptermodality{76}

\begin{chapteropening}
\textbf{By the end of this chapter:} \emph{I can form the plural of a word ending in a consonant, and say why the extra vowel is there.}

Complete ES-C44-plural-es: pluralise espanol, name the word that has carried this rule since long before the chapter, and say why -s alone will not do.
\end{chapteropening}

\section[-es]{-es — the plural your first chapter owed you}

\label{lesson:ES-C44-plural-es}



\begin{quote}
Say \textquotedblleft{}you,\textquotedblright{} formally. (\emph{Usted}.) Now say it to a group. (\emph{Ustedes}.) Hold both of those in your mouth for a moment.
\end{quote}

\begin{grammarlens}[title={Grammar Lens: you have been doing this for fifty chapters}]
Your first chapter gave you half a rule and said the rest was for later. This is later.

\begin{quote}
\textbf{A word ending in a vowel adds \emph{-s}. A word ending in a consonant adds \emph{-es}.}
\end{quote}

And you already know the second half, because you have been saying it since the day you learned to address a group:

\begin{quote}
\emph{usted} $\to$ \emph{usted\textbf{es}}.
\end{quote}

That is not a special word. That is the consonant rule, and nobody told you. Here it is again on something you also hold:

\begin{quote}
\emph{español} $\to$ \emph{español\textbf{es}} — Spanish people, Spanish speakers. \emph{los españoles} — with the plural article from last chapter.
\end{quote}
\end{grammarlens}

\begin{culture}
The extra \emph{e} is not decoration and not an exception. Try saying \emph{usteds} out loud. Then \emph{españols}.

Spanish does not let a word end in that pile of consonants. So when the plural \emph{-s} would land somewhere unsayable, the language puts in the vowel it needs to get the \emph{-s} out — and that vowel is \emph{e}.

Which means the rule is really \textbf{one} rule with a repair attached: \emph{add -s; if you cannot say it, make room}. That is a better thing to hold than two rules, because it will keep explaining things long after this chapter.
\end{culture}

\subsection*{Your turn}
Say these aloud:

\begin{itemize}
\raggedright
  \item \textquotedblleft{}usted\textquotedblright{} then \textquotedblleft{}ustedes\textquotedblright{} — the rule you already had
  \item \textquotedblleft{}los españoles\textquotedblright{}
  \item try \textquotedblleft{}usteds\textquotedblright{} and hear why the \emph{e} exists
\end{itemize}

\begin{tcolorbox}[breakable,colback=teal!4,colframe=teal!35!black,title={Before you move on}]
Vowel ending — which plural? (\textbf{-s}.) Consonant ending? (\textbf{-es}.) And why? (A bare \emph{-s} there would be \textbf{unsayable}.) Next: the four articles in one place, now that every one of them has been earned.
\end{tcolorbox}
